% ===================================================================
\chapter{Conclusions and Future Work}
\label{ch:conclusion}

\section{Summary}

This thesis set out to compare data compression methods for SoC-based radar on
real, public, multi-transmitter data, to select the best lossless ones, and to
implement their encoders and decoders on an FPGA and measure them on a board.
The four objectives of Section~\ref{sec:objectives} were met.

Eight methods were compared in MATLAB on 50 frames of ColoRadar data from a
twelve-transmitter, sixteen-receiver Texas Instruments cascade. Two lossless
predictive coders --- Doppler Redundancy Hybrid Encoding and a second-order
linear predictive coder --- were specified, implemented in C++ for Vitis HLS
with matching decompressors, and built with Vivado into a test system on an AMD
KCU116 board. They were verified bit for bit from MATLAB through C-simulation
and RTL co-simulation to the board, where the compressors' output matched the
simulation byte for byte and the hardware decoders reconstructed every frame
exactly.
A third, independent implementation in Python reproduced their bit counts.

\section{Findings}

\begin{enumerate}
\item \textbf{Only lossless predictive coding meets the requirement.} DRHE
(\num{3.75}) and LPC (\num{3.70}) are the only methods that compress by more
than three without adding a missed target. The lossy quantisers reach up to
\num{3.76} but miss \SIrange{10}{29}{\percent} of targets, and floating-point
storage compresses by only \num{1.04}. All detection loss in the lossless chain
comes from the 16-bit quantisation step, not from compression.

\item \textbf{On a TDM-MIMO radar, predict from the same transmitter.} Slow-time
prediction from the adjacent ramp, which on this radar is another transmitter,
compresses worse than no prediction at all. Predicting $\nTx = 12$ ramps back,
from the same antenna, makes prediction worth \SIrange{6}{8}{\percent}. For
DRHE the change is an index into the state and nothing else.

\item \textbf{Most of the ratio comes before prediction.} Huffman coding alone
compresses this data by \num{3.48}, mainly because the fixed 16-bit scaling
leaves six bits of every value unused and because a range-FFT frame is sparse.
Filling those bits removes the quantisation error but also almost all of the
compression. A no-prediction baseline should be reported with any radar
compression ratio.

\item \textbf{The dictionary matters more than the predictor.} A Huffman
dictionary re-derived from five frames for each algorithm's prediction errors
raises DRHE from \num{3.75} to \num{4.22} and LPC from \num{3.70} to \num{4.40}
on the board, holds on the 45 frames it was not trained on, changes only the
contents of a table in hardware, and puts LPC ahead of DRHE on ratio.

\item \textbf{The hardware is exact and fast.} On the KCU116 at
\SI{100}{\mega\hertz}, both compressors produced output byte-identical to the
simulation on all 50 frames, and the board's own decoders reconstructed every
frame exactly. DRHE
compresses a frame in \SI{0.61}{\milli\second} (\SI{5.1}{Gbit/s}) with
\SI{13}{\micro\second} of latency; LPC in \SI{3.25}{\milli\second} with
\SI{3.0}{\milli\second} of latency, because it must see the whole frame first.
The decoders take \SI{1.97}{\milli\second} and \SI{15.4}{\milli\second}; the
LPC decoder's time comes from a serial codeword search, not from the
algorithm.
Encoder and decoder together use \SI{23}{\percent} and \SI{25}{\percent} of the
chip's logic.

\item \textbf{DRHE is the better choice for a streaming sensor.} The two
algorithms are close on ratio and level on logic, but DRHE needs no frame
buffer, streams its output, and scales with the data rate, while LPC's buffer
and latency grow with the frame until, at the frame size of Kiem's use case, it
could not keep up.

\item \textbf{The baseline carries over, with one change.} Kiem's DRHE is
reproduced bit for bit and works on a cascade once the prediction lag is set to
the transmitter count; with a matched dictionary it compresses public cascade
data by \num{4.2}. Unlike his system it has a hardware decoder, and the
board's compressed output reconstructs exactly on the board.
\end{enumerate}

\section{Future work}
\label{sec:futurework}

\paragraph{Measure power, and speed up the LPC decoder.} Measure the board's
power draw for each build, to give the energy cost of compression. Replace the
LPC decoder's serial codeword search with the parallel matcher of the DRHE
decoder, which should cut its \SI{15.4}{\milli\second} per frame to about
\SI{2}{\milli\second}.

\paragraph{Move DRHE to fixed-point arithmetic.} The floating-point datapath is
the most expensive design decision and buys \SI{0.07}{\percent} of ratio.
Kiem's fixed-point types would remove most of its logic, DSPs and pipeline
depth while remaining exactly lossless, since the closed loop makes the decoder
follow whatever arithmetic the encoder uses. With fixed point, reading the state
as zero on the first ramps instead of writing it would also reach
$\mathrm{II} = 1$ almost for free.

\paragraph{A streaming LPC.} A backward-adaptive LPC that fits its coefficients
to previous frames would need no frame buffer and add no frame of latency,
combining LPC's better prediction with DRHE's streaming. For the forward LPC,
overlapping the three passes and amortising the per-bin pipeline fill would
cut its time per frame by about a third, and processing several bins in
parallel by more.

\paragraph{Dictionaries that adapt.} Train dictionaries on more varied scenes,
or let the hardware hold several and choose per frame, or update one slowly
from past frames. A seventeenth category for the value $-32768$ would remove
the need for saturation.

\paragraph{All sixteen channels and the range FFT on chip.} Widen the stream or
time-multiplex the cores to compress all sixteen receive channels, and place a
range-FFT core in front of the compressor, as in Kiem's system, so that the
board processes raw ADC data end to end.

\paragraph{The quantisation trade-off.} Study adaptive or companded scaling as
a deliberate design choice: how many bits of the noise floor are worth keeping
for the detection margin they buy, and what that costs in link bandwidth.

\paragraph{Other multiplexing schemes and methods.} Doppler- and code-division
MIMO distribute the slow-time structure differently, and whether slow-time
prediction helps on them is open. Learned and spectral methods such as AdaRadar
\cite{park2026} define the upper end of achievable ratios and could be compared
against the lossless coders on the same data.
